\section{Operator Theory}\label{chap:operator-theory}

In this chapter we give a brief discussion of the theory of linear operators $A$ from a Banach space $X$ to a Banach space $Y$. Our primary concerns center on the equation
\begin{equation}\label{eq:abstract-linear-operator-equation}
Ax=y,
\end{equation}
where $y\in Y$ is given, and the main issues we address are existence, multiplicity, and computability of solutions $x\in X$. Of course, most readers have already addressed these issues in studying linear algebra. There, the spaces $X$ and $Y$ are the finite-dimensional vector spaces $\mathbb{R}^n$ and $\mathbb{R}^m$, respectively, and $A$ is represented by an $m\times n$ matrix. We have already considered a more general type of operator in this text when we defined a bounded linear operator from one (possibly infinite-dimensional) Banach space to another in Definition~\ref{def:bounded-linear-mapping}. However, as we shall see below, many important operators in PDEs (and ODEs) are unbounded. The reader is strongly encouraged to compare the results of this section with the results of his or her old linear algebra text while keeping in mind the two main extensions of the theory: to spaces that are infinite-dimensional and to operators that are unbounded.

\textbf{Note:} Although we have defined operators to be maps between Banach spaces, most of the applications of operator theory that we address in this book will be to maps between separable Hilbert spaces. Thus, in many of the theorems below, we have given either statements or proofs only for the case of Hilbert spaces or separable Hilbert spaces. This practice greatly reduces the amount of machinery we need to develop, but it also limits the possible applications one can address using only material from this book. This is one of the prices you pay for learning functional analysis ``in the street.''

In the following, we will use the notations $X$ and $Y$ to refer to Banach spaces and $H$ to refer to a Hilbert space unless we specify otherwise.

\subsection{Basic Definitions and Examples}\label{sec:operator-basic-definitions-examples}

\subsubsection{Operators}\label{subsec:operators}

In order to accommodate unbounded operators we begin this section with the following extended definition.

\begin{Definition}
\label{def:linear-operator-with-specified-domain}
Let $X$ and $Y$ be Banach spaces. A linear operator from $X$ to $Y$ is a pair $(\mathcal{D}(A),A)$ consisting of a subspace $\mathcal{D}(A)\subset X$ (called the \textbf{domain} of the operator) and a linear transformation $A:\mathcal{D}(A)\to Y$.
\end{Definition}

Many mathematics students have had to endure a calculus teacher who insisted that there was a profound difference between the function $f(x)=x$ with domain $[0,1]$ and the same function defined on the whole real line. The students soon realize that in most cases the distinction can be ignored. In the course of this chapter, we shall see that including the domain in the definition of an operator is more than just pedantry. For unbounded operators, the specification of the domain can make a real difference. However, after having made such a big deal of the importance of the domain in the definition of an operator, we will often use sloppy language which ignores the point. That is, we will often refer to ``the operator $A$'' and leave the domain unspecified. This usage is standard and unambiguous in the study of bounded operators (whose domain, we see in Theorem~\ref{thm:bounded-operator-norm-preserving-extension} below, can be extended to all of $X$), and when there is no chance of confusion, we simply stick with the shorter nomenclature even for unbounded operators.

We will use both of the notations $Ax$ and $A(x)$ to indicate the action of an operator on elements of its domain.

\begin{Definition}
\label{def:operator-range-null-space}
The range of $(\mathcal{D}(A),A)$ is a subspace $\mathcal{R}(A)\subset Y$ defined by
\begin{equation}\label{eq:linear-operator-range-definition}
\mathcal{R}(A):=\{u\in Y\mid u=A(x),\quad\text{for some }x\in\mathcal{D}(A)\}.
\end{equation}
The null space of $(\mathcal{D}(A),A)$ is the subspace $\mathcal{N}(A)\subset X$ defined by
\begin{equation}\label{eq:linear-operator-null-space-definition}
\mathcal{N}(A):=\{x\in X\mid A(x)=0\}.
\end{equation}

\end{Definition}

With the range thus defined, we can use the following notation for the operator $(\mathcal{D}(A),A)$:\[X\supset\mathcal{D}(A)\ni x\mapsto A(x)\in\mathcal{R}(A)\subset Y.\]

The sets $X$ and $Y$ are sometimes referred to as the \emph{corange} and the \emph{codomain} in order to distinguish them from their subspaces, the domain and range, respectively. Although we agree with the importance of the distinction, we shall not adopt these terms.
